% Notes of the Montreal seminar, 1970 (cote 129, typed part): archive cote 129, site page 33 (PDF page 34).
% typescript. Draft compared: none.
% AI-transcribed by Claude Opus 5.5 (claude-opus-5-5, adaptive thinking, effort high), Message Batches API, batch mtl129_full (msgbatch_01AGnDKHwy6ATCM7ehUa6263), 2026-11-13; not yet checked by eye.
\scanpage{129/33}
\mspage{\hand{II.5}}

On prouve (SGA 3, V, 4.1) que $G''$ est \uline{représentable par un}
\margin{$\pi$} \uline{schéma en groupes finis localement libre sur} $S$ , et que \deleted{$p$}\inserted{$\pi$} est fidèle-%
ment plat, de sorte que $G$ est fini localement libre sur $G''$ , de rang
relatif égal au rang de $G'$ . Par transitivité, on en conclut donc que
\[ \mathrm{rang}(G) = \mathrm{rang}(G')\cdot\mathrm{rang}(G'') . \]

Par ailleurs, on montre que si dans la suite exacte de fais-%
ceaux ci-dessus, on suppose $G'$ et $G''$ représentables par des schémas
en groupes finis localement libres sur $S$ , alors $G$ est lui aussi repré-%
sentable par un groupe fini localement libre sur $S$ . C'est-à-dire qu'une
extension de groupes finis localement libres est encore finie localement
libre.

1.7. Soit $\ell$ un nombre premier. On appelle \uline{composante $\ell$-primaire}
du groupe $G$ sur $S$ le sous-groupe
\[ C(\ell) = \varinjlim_{n} \mathrm{Ker}\{G \xrightarrow{\ell^{n}\mathrm{Id}} G\} \quad ; \]
\margin{$\times$ $\ell$} c'est le plus grand sous-groupe de \hand{$\ell$}-torsion de $G$ .

Si $G$ est fini localement libre sur $S$ , on prouve que $G$ est
localement annulé par un entier $n$ (SGA 3, VIII, 7.3), et donc chacune de
ses composantes $\ell$-primaires est localement annulée par une puissance de $\ell$
\margin{$\times$ $\ell$} (i.e. est un \hand{$\ell$}-groupe), et $G$ admet alors une décomposition canonique
(comme faisceau f.p.p.f.) en ses composantes $\ell$-primaires :
\[ G = \coprod_{\ell} G(\ell) . \]

\margin{33}
